\section{Introduction}
\label{sec:intro}

互联网流量的异构化已成为骨干网面临的核心挑战。同一基础设施需同时承载交互式视频、远程控制、分布式 AI 训练与批量备份等业务，这些业务对带宽、时延、丢包等 QoS 维度的偏好差异巨大，形成了``老鼠流''（带宽小、QoS 要求严苛）与``大象流''（带宽巨大、容忍尽力而为）共存的混合特征\cite{CISCO2023_GlobalTraffic}。这种差异化需求使得传统无差异化的路由策略——无论是基于最短跳数的OSPF 还是等价多路径的 ECMP——在资源利用与 QoS 保障之间顾此失彼。

问题的本质困难在于在线因果耦合：流逐条到达且未来不可见，智能体为第$t$ 条流选路时无法预见 $t+1,\dots,W$ 的后续需求\cite{Wang1996_QoSRouting}。若允许离线处理整条流序列，选路可建模为整数规划求得全局最优；但在线约束迫使决策必须在信息不完备下做出，导致先到流与后到流竞争共享链路容量。贪心启发式（如 CSPF）虽每步局部最优，却常将早期到达的老鼠流优先分配至资源充裕的最优路径，耗尽后续大象流所需的链路容量，这正是短视启发式策略的经典困境\cite{Fortz2000_TrafficEngineering}。最优策略必须非短视地在需求分布的期望意义下预留容量、跨流权衡，并将序列末尾的累积效用回溯归因到早先各步——这种``未来不可见下的序贯决策与长程信用分配''的结构，天然适合用强化学习（RL）求解。

近年来，深度强化学习（DRL）与图神经网络（GNN）的结合为在线路由带来了新的可能：DRL 提供了非短视的序贯决策能力\cite{Sun2022_ScalableRouting_TON,He2024_RoutingOptimization_TMC}，GNN 提供了对网络拓扑的显式编码能力\cite{Rusek2020_JSAC_RouteNet, FerriolGalmes_2023_RouteNetFermi}。然而，现有方法存在两个深层局限：其一，DRL 路由策略通常为所有流优化单一全局目标，无法在混合业务下实现差异化服务；其二，GNN 方法为所有流学习统一的拓扑嵌入，无法根据逐流的 QoS 特征对同一拓扑进行差异化感知，即``同图同嵌入''问题。

要打破``同图同嵌入''的同质化限制，需要使网络编码器具备流条件化能力，即根据每条流的个体 QoS 特征动态调整其拓扑感知。条件化神经网络（conditional neural networks）为此提供了核心工具。Perez 等人提出的 FiLM（Feature-wise Linear Modulation）通过仿射变换将辅助信息注入网络各层\cite{Perez2018_FiLM}，Brockschmidt 等人将其拓展至图领域（GNN-FiLM），以目标节点表示调制消息传递\cite{Brockschmidt2020_GNNFiLM}。但在在线路由场景中，真正需要条件化的并非目标节点，而是逐流的 QoS 意图：同一张拓扑在面对老鼠流与大象流时，其消息传递过程应当被重塑为两种截然不同的拓扑感知。现有文献尚未将 FiLM 条件化机制引入逐流在线路由场景，使得混合业务下的差异化拓扑感知与策略学习成为开放问题。

针对上述差距，本文提出流感知 FiLM-GNN 编码器，核心思想是：通过 FiLM 将每条流的 QoS 意图向量注入 GNN 的每一层消息传递，使同一张物理拓扑在不同意图下被编码为差异化的节点嵌入。进而，我们将在线混合业务选路形式化为马尔可夫决策过程（MDP），以近端策略优化（PPO）训练非短视选路策略。本文的主要贡献如下：
\begin{itemize}
    \item \textbf{面向混合业务在线路由的流感知图强化学习架构：}
    该架构由流条件化图编码器、路径感知打分器与独立价值网络三部分构成，通过近端策略优化实现端到端训练，解决了同一拓扑需对不同QoS意图流产生差异化路由决策的核心矛盾。

    \item \textbf{逐流QoS意图驱动的FiLM条件化机制与网络设计：}
    以逐流意图向量生成图神经网络各层的仿射调制参数，使消息传递过程成为流感知的；配套设计了基于LSTM的路径嵌入池化、意图-路径联合打分及意图条件化注意力图池化等关键模块，实现了从拓扑编码到路径选择的完整流条件化决策链路。

    \item \textbf{双空间可解释性实验验证：}
    实验表明所提架构在累积效用、带宽满足比和端到端时延上显著优于经典基线；进一步通过表示空间（隐藏特征分离度）与决策空间（路径选择分化）的双重分析证实：仅改变QoS意图即可在相同拓扑、相同源宿节点对下触发差异化路由决策，从机理层面验证了FiLM条件化的必要性。
\end{itemize}
